\section{Before any task}
\label{sec:before}

\subsection{How the tasks depend on one another}

The tasks are not independent. Photographs can be taken only after the board works and the
camera stands at a known height (Task~1). The classifier is retrained on real photographs of
the symbol sheets (Task~2) before the pipeline is tuned (Task~3), because tuning a pipeline
whose classifier is about to change wastes effort. Tuning requires photographs of the 20
development lines, which belong to the printed material of Task~4; and the final accuracy
figures of Task~4 may be measured only after tuning has stopped, since a measurement that has
been used for tuning no longer measures anything. Task~4 is therefore carried out in two parts:
its photography early, its measurement late. Figure~\ref{fig:order} shows the resulting order,
and Table~\ref{tab:order} the days on which it fits the dates of the update report.

\begin{figure}[H]
\centering
\begin{tikzpicture}[node distance=5mm and 6mm, font=\small,
  box/.style={draw, rounded corners=2pt, align=center, minimum height=11mm, inner sep=3pt, text width=24mm},
  >={Stealth[length=2mm]}]
\node[box] (t1) {Task 1\\\scriptsize board, stand};
\node[box, right=of t1] (t2) {Task 2\\\scriptsize sheets, retraining};
\node[box, right=of t2] (t4a) {Task 4, part A\\\scriptsize photograph 80 lines};
\node[box, right=of t4a] (t3) {Task 3\\\scriptsize triage, freeze};
\node[box, below=9mm of t3] (t4b) {Task 4, part B\\\scriptsize final accuracy};
\node[box, left=of t4b] (t5) {Task 5\\\scriptsize demo video};
\node[box, left=of t5] (t6) {Task 6\\\scriptsize report, rehearsal};
\draw[->] (t1) -- (t2);
\draw[->] (t2) -- (t4a);
\draw[->] (t4a) -- (t3);
\draw[->] (t3) -- (t4b);
\draw[->] (t4b) -- (t5);
\draw[->] (t5) -- (t6);
\end{tikzpicture}
\caption{Recommended order of work. The photography of Task~4, part~A, fits into the hour
during which the classifier of Task~2 trains.}
\label{fig:order}
\end{figure}

\begin{table}[H]
\centering
\begin{tabular}{C{2.5cm}L{6.2cm}C{2.4cm}}
\toprule
\hd{Day} & \multicolumn{1}{c}{\hd{Work}} & \hd{Time needed} \\
\midrule
Tue Nov 24 & Task 1: board, soak test, stand, calibration & 3 h \\
Wed Nov 25 & Task 2: photograph the sheets, crop, start the training & 1.5 h \\
Wed Nov 25 & Task 4, part A: photograph the 80 lines while the classifier trains & 1 h \\
Wed Nov 25 & Task 2, continued: evaluate and upload the classifier & 0.5 h \\
Wed Nov 25 & Task 3: triage, trials, decision on fractions, freeze & 2 h \\
by Mon Nov 30 & Task 4, part B: the final measurement & 0.5 h \\
by Mon Nov 30 & Task 5: demo video & 1.5 h \\
Mon Nov 30 & Task 6: two rehearsals & 1 h \\
Tue Dec 1 & Task 6: report, presentation, release, hand-in & as required \\
\bottomrule
\end{tabular}
\caption{A schedule that satisfies the dates of the update report. Nothing is planned for the
Thanksgiving holiday, November 26 to 29.}
\label{tab:order}
\end{table}

\subsection{Conventions}

This document follows the conventions of the Set-up Guide. A blue box contains one command,
which is typed or pasted as a single line and completed with \key{Enter}. A gray box shows
what the screen should then display. A step ends with a \ok~\textbf{Checkpoint}, an observable
fact that shows the step has succeeded; do not continue until it is true. An amber box says
what to do when it is not. Words to be replaced by your own are printed in capitals on a
yellow background; this document contains a single such placeholder, \ph{NAMES}, which stands
for a list of photograph names and is explained where it occurs. Where the output of a program
depends on a measurement that has not been made yet, the gray box prints three dots in place
of the figures.

\subsection{Conditions that hold for every task}

\step{Confirm that the laptop is ready}\label{st:ready}
Open a new PowerShell window (click \ui{Start}, type \code|powershell|, click \ui{Windows
PowerShell}) and give the opening sequence of the Set-up Guide.
\begin{ps}[Opening sequence, command 1 of 2]
cd C:\ee105\camera-to-latex
\end{ps}
\begin{ps}[Opening sequence, command 2 of 2]
.\.venv\Scripts\Activate.ps1
\end{ps}
\checkpoint{The prompt reads:}
\venvprompt
If your copy of the project is in another folder (the project lead's is), use the path of that
folder in the first command, here and throughout this document.

\step{Bring the project up to date}\label{st:uptodate}
\begin{ps}
git pull
\end{ps}
\begin{ps}
python -m pytest -q
\end{ps}
\begin{out}
................................................                         [100%]
48 passed in 9.87s
\end{out}
\checkpoint{The last line says \emph{48 passed} (the time will differ, and a remark about
warnings is harmless). If it says \emph{45 passed}, your copy does not yet contain the
corrections of 24 November (the camera-model setting of the sketch and the more tolerant sheet
cropper), on which Tasks~1 and~2 rely: ask the lead whether they have been uploaded, and repeat
this step.}

\step{Confirm that a classifier is present}
\begin{ps}
Get-ChildItem models
\end{ps}
\checkpoint{The listing contains two files whose names begin with \file{symbols\_v1}; one of
them ends in \file{.keras} and the other in \file{.json}. If the folder does not exist,
complete the quick start of the Set-up Guide first, or fetch the team's classifier with
\cmd{git pull} once Task~2 has published it.}

\subsection{How to change a setting}
\label{sec:setting}

Every adjustable number of the pipeline is a line of the file \file{config.py} of the form
\emph{NAME = value}, followed by a remark that begins with the sign \code|#|. Several tasks
change one of these lines. The method is always the same.

\step{Open the file}
\begin{ps}
notepad config.py
\end{ps}
Notepad opens with the file. Press \key{Ctrl}+\key{F}, type the name of the setting exactly as
this document prints it, and press \key{Enter}; Notepad jumps to the line.

\step{Change the value, and nothing else}
Replace only the value between the equals sign and the \code|#| sign. Leave the name, the
equals sign, the spacing at the start of the line and the remark as they are; a line of Python
that begins with an extra space is an error. Capital letters matter: the two logical values are
written \code|True| and \code|False|.
\begin{filebox}{Example: the line FLIP\_H of config.py, before and after}
FLIP_H = False
FLIP_H = True
\end{filebox}
Press \key{Ctrl}+\key{S} to save, and close Notepad.

\step{Check that the file is still valid}\label{st:valid}
\begin{ps}
python -c "import config; print('config.py is valid')"
\end{ps}
\checkpoint{The answer is \emph{config.py is valid}. Any other answer names the line that is
damaged; open the file again and compare that line with the examples in this document. To
discard all changes made since the last commit and start again, run
\texttt{git restore config.py}.}

\subsection{How results are shared}

Each task ends by recording its results in the repository, with the three commands \cmd{git
add}, \cmd{git commit} and \cmd{git push} written out for that task. Section~9 of the Set-up
Guide explains them and says what to do when Git objects. Photographs of the symbol sheets and
other large working files are not put into the repository; where they must be kept, the task
says where.

\subsection{Handling the board}

\begin{caution}[Read before touching the board]
\begin{itemize}
\item The board is an unprotected circuit. Keep metal objects (metal rulers, clips, jewellery,
  the spiral of a notebook) and liquids away from it, and never lay it on a conductive surface
  while the USB cable is connected.
\item Hold it by its edges. Do not press on the camera module, on its lens, or on the Arduino
  Nano.
\item Insert and remove the USB plug straight, holding the Nano's edge with the other hand, and
  fasten the cable to the stand so that a pull on the cable cannot move the board.
\item The board belongs to the course. Note its number (printed on the white disc beside the
  camera) in the record of Task~1, and store it in its box after each session.
\end{itemize}
\end{caution}

\subsection{How precise a measured percentage is}
\label{sec:stat}

Every accuracy in this project is a share $p$ of correct results among $n$ trials: 20
development lines, 60 test lines, 280 held-out crops. Such a share is itself uncertain, because
another set of $n$ lines of the same kind would give a somewhat different value. For
independent trials the size of this uncertainty, the \term{standard error}, is
\[
  \sigma_p = \sqrt{\frac{p\,(1-p)}{n}} ,
\]
and the true accuracy lies within about two standard errors of the measured share in nineteen
cases out of twenty. Table~\ref{tab:stat} evaluates the formula for the sample sizes of this
project.

\begin{table}[H]
\centering
\begin{tabular}{L{5.3cm}C{1.5cm}C{1.9cm}C{2.2cm}}
\toprule
\multicolumn{1}{c}{\hd{Measurement}} & \hd{Trials $n$} & \hd{Share $p$} & \hd{Standard error} \\
\midrule
Fraction lines of the test set & 10 & 60 \% & 15 points \\
Development lines & 20 & 70 \% & 10 points \\
Script lines of the test set & 20 & 70 \% & 10 points \\
Plain lines of the test set & 25 & 80 \% & 8 points \\
All test lines & 60 & 75 \% & 6 points \\
Simulated lines, two sets & 160 & 80 \% & 3 points \\
Held-out real crops & 280 & 95 \% & 1.3 points \\
\bottomrule
\end{tabular}
\caption{Standard error of a share for the sample sizes of this project.}
\label{tab:stat}
\end{table}

Two practical rules follow. First, a result is always recorded as a count, such as 17 of 25,
and not only as a percentage, so that the reader can judge its precision. Second, a difference
of one or two lines between two settings, measured on 20 lines, proves nothing by itself;
Task~3 therefore accepts a change only when it helps on the real photographs and does no harm
on a larger set of simulated ones.

\subsection{Records}

Each task section ends with a short table to be filled in. These records are the raw material of
the final report (Task~6): the report must state the camera height, the number of pixels per
centimetre, the accuracies and the dates, and it is far easier to copy them from a record than
to reconstruct them. Fill the tables in on a printed copy of this document or copy them into a
shared file; Appendix~\ref{sec:records} lists them with their pages.

\begin{note}[Status of this document]
The procedures that involve only the laptop (cropping, retraining, evaluation, review and the
release commands) were executed in a test environment on 24 November 2026 with simulated
photographs, and the quoted outputs come from those runs or from the project lead's runs of 23
November. The procedures that involve the board, the camera and the stand have not been
executed: they are the tasks themselves. Wherever a quantity depends on the hardware (field of
view, sharpness, exposure), the procedure therefore measures it and decides from the
measurement, and the numbers given in advance are predictions.
\end{note}
